% Generated draft: requires mathematical and source review.
\documentclass[11pt]{article}
\usepackage[margin=25mm]{geometry}
\usepackage{fontspec}
\setmainfont{Latin Modern Roman}
\usepackage{amsmath,amssymb,mathtools,unicode-math}
\setmathfont{Latin Modern Math}
\usepackage{enumitem,needspace,xurl,hyperref}
\setlength{\parindent}{0pt}
\setlength{\parskip}{5pt}
\setlength{\emergencystretch}{3em}
\newcommand{\N}{\mathbb{N}}
\newcommand{\Z}{\mathbb{Z}}
\newcommand{\Q}{\mathbb{Q}}
\newcommand{\R}{\mathbb{R}}
\newcommand{\C}{\mathbb{C}}
\newcommand{\F}{\mathbb{F}}
\newcommand{\A}{\mathbb{A}}
\newcommand{\PP}{\mathbb{P}}
\newcommand{\E}{\mathbb{E}}
\newcommand{\eps}{\varepsilon}
\newcommand{\dd}{\,\mathrm{d}}
\begin{document}
\section*{UnsolvedMath ID: 30004849}
\textbf{Source problem code:} catalogue entry
\subsection{Definitions and mathematical statement}
A two-dimensional topological conformal field theory (2d--TCFT) and a two-dimensional homological conformal field theory (2d--HCFT) are the two kinds of field theory addressed by this record. The supplied original excerpt indicates a chain-level cobordism formulation for the former and a symmetric monoidal functor to graded vector spaces for the latter, but its notation is corrupted and its source categories are not defined.

To display the intended functorial framework without supplying missing hypotheses, suppose a coefficient field \(k\) and appropriate symmetric monoidal source categories \(\mathcal B_{\mathrm{ch}}\) and \(\mathcal B_{\mathrm H}\) have been specified. Write \(\mathsf{Ch}(k)\) for the differential graded category of chain complexes of \(k\)-vector spaces and \(\mathsf{gVect}_k\) for the category of graded \(k\)-vector spaces. The fragment suggests schematic functorial expressions
\[
F_{\mathrm{ch}}\colon \mathcal B_{\mathrm{ch}}\longrightarrow\mathsf{Ch}(k),
\qquad
F_{\mathrm H}\colon \mathcal B_{\mathrm H}\longrightarrow\mathsf{gVect}_k.
\]
Here the source categories are placeholders, not recovered definitions of the damaged cobordism notation. A symmetric monoidal functor preserves the tensor product and its unit through coherent structure isomorphisms compatible with symmetry. The excerpt explicitly describes the HCFT functor as symmetric monoidal; the precise enrichment and coherence requirements for the intended theories require clarification from the source. The coefficient field, grading conventions, and tensor conventions are not supplied, so the displayed framework is conditional and does not select them.

\textbf{Research questions.} For each of the intended notions of 2d--TCFT and 2d--HCFT:
\begin{enumerate}
\item What algebraic structures describe these theories, and what classification statement provides the analogue of the theorem referred to in the original discussion?
\item Do nontrivial examples exist?
\item What mathematical applications do such theories have?
\end{enumerate}
The referenced theorem is absent from the supplied record. The intended cobordism categories, permitted boundary conditions, meaning of nontriviality, and notion of equivalence needed for a classification are also unspecified. Consequently these are research questions with unresolved formulation data, rather than a fully determined conjecture with fixed hypotheses and quantifiers. The record does not justify choosing an open, closed, or open--closed version on its behalf.

The supplied dated assessment reports substantial classification results for open and open--closed TCFTs using Calabi--Yau \(A_\infty\)-categories, together with known nontrivial examples and applications, while describing the broader TCFT/HCFT questions as incomplete. This is supplied curation, not an independently verified mathematical or literature assessment.
\subsection{Short English statement}
Identify algebraic descriptions of two-dimensional topological conformal and homological conformal field theories, find nontrivial examples, and explain their applications. The precise field-theory categories and the earlier classification theorem are missing from the supplied excerpt. The supplied partially solved status therefore accompanies an incompletely specified research programme.
\subsection{Sources}
\begin{itemize}
\item Local UnsolvedMath registry, supplied record with permanent ID \(30004849\), exact source problem code \texttt{catalogue entry}, title ``Algebraic Structures and Applications of Two-Dimensional Field Theories,'' in the Oberwolfach Reports Open Problems collection. No individual proposer is identified; the supplied proposed year is 2021.
\item Supplied source citation: \emph{Homotopical Algebra and Higher Structures} (2022), Oberwolfach Reports, pp.\ 2497--2560, DOI \texttt{10.4171/owr/2021/46}. \url{https://doi.org/10.4171/owr/2021/46}. The background attributes the questions to workshop open-problem discussions documented in Oberwolfach Reports.
\item Supplied literature reference: Costello, ``TCFTs and Calabi--Yau categories.'' \url{https://doi.org/10.1016/j.aim.2006.06.004}.
\item Supplied literature reference: ``OWR report stating the remaining questions.'' \url{https://ems.press/content/serial-article-files/46922}.
\item Literature assessment dated exactly \texttt{2026-08-22}: the supplied curation labels the problem partially solved, reports a major algebraic classification of open and open--closed TCFTs through Calabi--Yau \(A_\infty\)-categories and known examples and applications, and regards the broader classification and usefulness questions as incomplete. This is supplied dated curation derived from the OWR record, not independently verified here; its time-sensitive status and cited claims require independent checking.
\item The registry labels the statement \texttt{corrected\_verified} and reports a manual semantic, typographic, and LaTeX review. These are supplied editorial assertions and do not resolve the missing mathematical definitions or constitute mathematical review of this draft.
\end{itemize}
\end{document}
